\section{Additional Notes on Benchmark Mapping}

The benchmark comparison in the main paper is intentionally qualitative. Kim et al.~\cite{kim2025scaling}
report architecture-level aggregates across heterogeneous tasks, models, and coordination regimes. Those
metrics are useful for checking whether the theory predicts the right ordering across topology classes,
but they should not be interpreted as direct fitted values of variables such as $d$, $\Delta I_j$, or
$S$ in the simplified formulas of Section~5.

This distinction matters especially for rebuttal. The paper does not claim that the theory numerically
reconstructs the benchmark tables. It claims that the benchmark outcomes are directionally consistent with
the visibility-based predictions: review bottlenecks reduce error amplification, artificial handoffs hurt
strictly sequential tasks, decomposition helps on approximately independent tasks, and large toolsets
increase coordination tax in distributed regimes.
